\documentclass[10pt,a4paper]{amsart}
\usepackage{amsmath,amssymb,amsthm,graphicx}
\usepackage{verbatim}
\newenvironment{reference}{\comment}{\endcomment}
%\newenvironment{reference}{}{}
\newenvironment{slogan}{\comment}{\endcomment}
\newenvironment{history}{\comment}{\endcomment}

% For commutative diagrams we use Xy-pic
\usepackage[all]{xy}

% We use 2cell for 2-commutative diagrams.
\xyoption{2cell}
% Theorem environments.
%
\theoremstyle{plain}
\newtheorem{theorem}[subsection]{Theorem}
\newtheorem{proposition}[subsection]{Proposition}
\newtheorem{lemma}[subsection]{Lemma}

\theoremstyle{definition}
\newtheorem{definition}[subsection]{Definition}
\newtheorem{example}[subsection]{Example}
\newtheorem{exercise}[subsection]{Exercise}
\newtheorem{situation}[subsection]{Situation}

\theoremstyle{remark}
\newtheorem{remark}[subsection]{Remark}
\newtheorem{remarks}[subsection]{Remarks}

\numberwithin{equation}{subsection}

% Macros
%
\def\lim{\mathop{\mathrm{lim}}\nolimits}
\def\colim{\mathop{\mathrm{colim}}\nolimits}
\def\Spec{\mathop{\mathrm{Spec}}}
\def\Hom{\mathop{\mathrm{Hom}}\nolimits}
\def\Ext{\mathop{\mathrm{Ext}}\nolimits}
\def\SheafHom{\mathop{\mathcal{H}\!\mathit{om}}\nolimits}
\def\SheafExt{\mathop{\mathcal{E}\!\mathit{xt}}\nolimits}
\def\Sch{\mathit{Sch}}
\def\Mor{\mathop{\mathrm{Mor}}\nolimits}
\def\Ob{\mathop{\mathrm{Ob}}\nolimits}
\def\Sh{\mathop{\mathit{Sh}}\nolimits}
\def\NL{\mathop{N\!L}\nolimits}
\def\CH{\mathop{\mathrm{CH}}\nolimits}
\def\proetale{{pro\text{-}\acute{e}tale}}
\def\etale{{\acute{e}tale}}
\def\QCoh{\mathit{QCoh}}
\def\Ker{\mathop{\mathrm{Ker}}}
\def\Im{\mathop{\mathrm{Im}}}
\def\Coker{\mathop{\mathrm{Coker}}}
\def\Coim{\mathop{\mathrm{Coim}}}

% Boxtimes
%
\DeclareMathSymbol{\boxtimes}{\mathbin}{AMSa}{"02}

%
% Macros for moduli stacks/spaces
%
\def\QCohstack{\mathcal{QC}\!\mathit{oh}}
\def\Cohstack{\mathcal{C}\!\mathit{oh}}
\def\Spacesstack{\mathcal{S}\!\mathit{paces}}
\def\Quotfunctor{\mathrm{Quot}}
\def\Hilbfunctor{\mathrm{Hilb}}
\def\Curvesstack{\mathcal{C}\!\mathit{urves}}
\def\Polarizedstack{\mathcal{P}\!\mathit{olarized}}
\def\Complexesstack{\mathcal{C}\!\mathit{omplexes}}
% \Pic is the operator that assigns to X its picard group, usage \Pic(X)
% \Picardstack_{X/B} denotes the Picard stack of X over B
% \Picardfunctor_{X/B} denotes the Picard functor of X over B
\def\Pic{\mathop{\mathrm{Pic}}\nolimits}
\def\Picardstack{\mathcal{P}\!\mathit{ic}}
\def\Picardfunctor{\mathrm{Pic}}
\def\Deformationcategory{\mathcal{D}\!\mathit{ef}}
\usepackage{cleveref}
\crefname{originalreference}{original reference}{original references}
\Crefname{originalreference}{Original reference}{Original references}
\makeatletter
\newlabel{lemma-construct}{{097C}{source}{}{sourcepointer-1}{}}
\newlabel{lemma-construct@cref}{{[originalreference][0][]097C}{[1][0][]source}{}{}{}}
\makeatother
\crefname{originalreference}{original reference}{original references}
\Crefname{originalreference}{Original reference}{Original references}
\makeatletter
\newlabel{lemma-ind-zariski-implies}{{096T}{source}{}{sourcepointer-2}{}}
\newlabel{lemma-ind-zariski-implies@cref}{{[originalreference][0][]096T}{[1][0][]source}{}{}{}}
\makeatother
\crefname{originalreference}{original reference}{original references}
\Crefname{originalreference}{Original reference}{Original references}
\makeatletter
\newlabel{lemma-local-isomorphism-fully-faithful}{{096L}{source}{}{sourcepointer-3}{}}
\newlabel{lemma-local-isomorphism-fully-faithful@cref}{{[originalreference][0][]096L}{[1][0][]source}{}{}{}}
\makeatother
\crefname{originalreference}{original reference}{original references}
\Crefname{originalreference}{Original reference}{Original references}
\makeatletter
\newlabel{lemma-silly}{{00DS}{source}{}{sourcepointer-4}{}}
\newlabel{lemma-silly@cref}{{[originalreference][0][]00DS}{[1][0][]source}{}{}{}}
\makeatother
\crefname{originallemma}{original lemma}{original lemmas}
\Crefname{originallemma}{Original lemma}{Original lemmas}
\makeatletter
\newlabel{topology-lemma-profinite}{{08ZY}{source}{}{sourcepointer-5}{}}
\newlabel{topology-lemma-profinite@cref}{{[originallemma][0][]08ZY}{[1][0][]source}{}{}{}}
\makeatother
\title{Realizing a profinite space as the connected components of an ind-Zariski cover}
\author{The Stacks Project Authors}
\date{stacks-project:proetale.tex:local-snapshot-20261012; GNU FDL 1.2 or later}
\begin{document}
\maketitle
\noindent\textit{Editorial note.} All delivered mathematical text is the author's own, copied from the pinned Stacks Project chapter file proetale.tex; no statement or proof step was rewritten, reordered or removed.\par\medskip
\noindent\textit{Editorial note.} The counted proof is the author's complete argument: the presentation of the profinite space as a limit of finite discrete spaces, the construction of the ind-Zariski quotients realizing the closed subsets of the components, and the passage to the colimit with the identification of the components.\par\medskip
\noindent\textit{Editorial note.} The label lemma-construct-profinite names the counted statement, and the proof environment immediately following it in the source is delivered verbatim as its proof.\par\medskip
\noindent\textit{Editorial note.} The counted statement is carried by the author's own lemma environment in the chapter source.\par\medskip
\noindent\textit{Editorial note.} Pre-open gate: python operations/agent-authoring/tools/precard\_gate.py --source proetale.tex --label lemma-construct-profinite --start 1097 --end 1139 --local-id 3593 returned RESULT=OK (L3593 free, occupancy FREE, public interval clear). Candidate taken from sources/stacks-proetale/candidates-halfA.json.\par\medskip
\noindent\textit{Editorial note.} The statement and proof are copied from the pinned chapter file; the cited inputs are kept as external pointers with their canonical Stacks tags.\par\medskip
\noindent\textit{Allowed standard prerequisites (editorial).}\begin{itemize}
\item Description of a profinite space as a limit of finite discrete spaces. Topology, the lemma describing a profinite space as the limit of an inverse system of finite discrete spaces. Supplies the presentation of the profinite space T as a limit of finite discrete spaces T\_i over a directed set.
\item Construction of an ind-Zariski ring map from a closed subset of the components. Same chapter, the lemma constructing an ind-Zariski quotient corresponding to a closed subset of the connected components. Supplies, for each i, the surjection A\_i -> B\_i which is ind-Zariski and realizes the closed subset Z\_i.
\item Products of rings and local isomorphisms. Same chapter, the lemma that a local isomorphism is fully faithful and the lemma producing ind-Zariski maps between the pieces. Supply the description of X x T\_i as the spectrum of the product of copies of A and the transition maps B\_\{i'\} -> B\_i.
\item Identification of the connected components of the limit. Same chapter, the lemma identifying the connected components of the limit with the limit of the components. Supplies the conclusion that T is the space of connected components of Y.
\end{itemize}\medskip
\ifdefined\hypertarget\hypertarget{sourcepointer-1}{}\fi
\noindent\textit{Source reference.} Original reference 097C: sources/stacks-proetale/pinned-original-source.tex, source lines 1081--1081. Prerequisite pointer invoked at proof line(s) 1118; it names a standard fact of the same chapter that the counted argument uses, and every step of the counted proof built on it is retained verbatim in the delivered body.\par\medskip
\ifdefined\hypertarget\hypertarget{sourcepointer-2}{}\fi
\noindent\textit{Source reference.} Original reference 096T: sources/stacks-proetale/pinned-original-source.tex, source lines 685--685. Prerequisite pointer invoked at proof line(s) 1126; it names a standard fact of the same chapter that the counted argument uses, and every step of the counted proof built on it is retained verbatim in the delivered body.\par\medskip
\ifdefined\hypertarget\hypertarget{sourcepointer-3}{}\fi
\noindent\textit{Source reference.} Original reference 096L: sources/stacks-proetale/pinned-original-source.tex, source lines 603--603. Prerequisite pointer invoked at proof line(s) 1125; it names a standard fact of the same chapter that the counted argument uses, and every step of the counted proof built on it is retained verbatim in the delivered body.\par\medskip
\ifdefined\hypertarget\hypertarget{sourcepointer-4}{}\fi
\noindent\textit{Source reference.} Original reference 00DS: sources/stacks-proetale/pinned-original-source.tex, source lines 434--434. Prerequisite pointer invoked at proof line(s) 1137; it names a standard fact of the same chapter that the counted argument uses, and every step of the counted proof built on it is retained verbatim in the delivered body.\par\medskip
\ifdefined\hypertarget\hypertarget{sourcepointer-5}{}\fi
\noindent\textit{Source reference.} Original lemma 08ZY: sources/stacks-chow/tags.txt, source lines 11391--11391. Prerequisite pointer invoked at proof line(s) 1115; it names a standard fact of a sibling chapter that the counted argument uses, and every step of the counted proof built on it is retained verbatim in the delivered body.\par\medskip
\begin{lemma}
\label{lemma-construct-profinite}
Let $A$ be a ring and let $X = \Spec(A)$. Let $T$ be a profinite space and
let $T \to \pi_0(X)$ be a continuous map. There exists an
ind-Zariski ring map $A \to B$ such that with $Y = \Spec(B)$ the diagram
$$
\xymatrix{
Y \ar[r] \ar[d] & \pi_0(Y) \ar[d] \\
X \ar[r] & \pi_0(X)
}
$$
is cartesian in the category of topological spaces and such that
$\pi_0(Y) = T$ as spaces over $\pi_0(X)$.
\end{lemma}

\begin{proof}
Namely, write $T = \lim T_i$ as the limit of an inverse system finite
discrete spaces over a directed set (see
Topology, Lemma \ref{topology-lemma-profinite}). For each $i$ let
$Z_i = \Im(T \to \pi_0(X) \times T_i)$. This is a closed subset.
Observe that $X \times T_i$ is the spectrum of $A_i = \prod_{t \in T_i} A$
and that $A \to A_i$ is a local isomorphism. By Lemma \ref{lemma-construct}
we see that $Z_i \subset \pi_0(X \times T_i) = \pi_0(X) \times T_i$
corresponds to a surjection $A_i \to B_i$ which is ind-Zariski
such that $\Spec(B_i) = X \times_{\pi_0(X)} Z_i$ as subsets of
$X \times T_i$. The transition maps $T_i \to T_{i'}$ induce maps
$Z_i \to Z_{i'}$ and $X \times_{\pi_0(X)} Z_i \to X \times_{\pi_0(X)} Z_{i'}$.
Hence ring maps $B_{i'} \to B_i$
(Lemmas \ref{lemma-local-isomorphism-fully-faithful} and
\ref{lemma-ind-zariski-implies}).
Set $B = \colim B_i$. Because $T = \lim Z_i$ we have
$X \times_{\pi_0(X)} T = \lim  X \times_{\pi_0(X)} Z_i$
and hence $Y = \Spec(B) = \lim \Spec(B_i)$
fits into the cartesian diagram
$$
\xymatrix{
Y \ar[r] \ar[d] & T \ar[d] \\
X \ar[r] & \pi_0(X)
}
$$
of topological spaces. By Lemma \ref{lemma-silly}
we conclude that $T = \pi_0(Y)$.
\end{proof}

\end{document}
